\section{Introduction}\label{sec:intro}

Online political conversations on social media increasingly serve as a vector for polarization and antagonism between groups with opposing political viewpoints. Platform ranking algorithms can amplify sensational and misleading content \cite{Strayetal2023AlgorithmicManagement}, and emergent communities form around shared grievances or shared antagonisms. A growing body of research has sought to leverage AI and natural-language technologies to address this---from LLM dialogues to reduce conspiracy beliefs \cite{Costelloetal2024ReducingConspiracyBeliefs} to consensus-finding tools that aggregate political opinions \cite{Tessleretal2024HabermasMachine, Argyleetal2023LeveragingAI}, to crowdsourced fact-checking systems augmented by language models \cite{Deetal2024Supernotes, LiAndBakker2026AIFactChecking}.

Most of these interventions, however, implicitly operate within a \textit{deliberative} model of democracy in which the goal is to dampen affect, secure rational consensus, and treat passionate disagreement as a defect to be repaired \cite{LazarAndManuali2026LLMsDemocracy}. We take a different starting point. Following Mouffe's notion of \textit{agonistic pluralism} \cite{Mouffe2000DemocraticParadox}, we ask what it would mean for an AI intervention to transform \textit{antagonistic} political relationships — in which one's opponent is an enemy to be destroyed — into \textit{agonistic} ones — in which one's opponent is an adversary to be contested, but whose right to a position is not put into question. On this view, affective polarization need not be uniformly suppressed; what matters is whether channels exist to convert antagonistic affect into agonistic forms \cite{Westphal2024AffectivePolarization}. This reframing has direct consequences for design: rather than asking how to design an intervention to maximise consensus, we ask how could one encourage disagreement while defusing enemy framings.

Drawing on prior work suggesting that humor can suspend conflict and recast political opponents as ``legitimate discussion partners'' \cite{BoukesAndHameleers2023FactsOrHumor, NorrickAndSpitz2008HumorAsAResource}, we we examine whether a satirical AI fact-checker can de-antagonize political conversations.

We operationalise this through Eddie Bot, a fact-checking AI reply bot deployed on X. When mentioned by an `invoker' on a post they wish to contest, Eddie Bot retrieves live web evidence against a curated source-quality table, and replies with one of five actions: \textit{verify}, \textit{provide context}, \textit{challenge opinion}, \textit{surface perspectives}, or \textit{decline}. Each reply is rendered in one of three tone registers — \textit{agreeable}, \textit{neutral}, or \textit{satirical.} We randomly assign a tone to each participant at enrollment. 

While our study is still in pilot, we aim deploy Eddie Bot in a between-subjects randomized controlled trial with 30 invokers over five days. We measure participants' self-report affective polarization (using scales adapted from \citet{Strayetal2026ProsocialRrankingChallenge} and \citet{Westphal2024AffectivePolarization}) and bystander engagement (likes, replies, link clicks, and LLM-based classifiers for antagonistic affective polarization in third-party replies). Our contributions are threefold: \textbf{(1)} a theoretical framing for AI interventions in online political conversations grounded in agonistic rather than deliberative democracy; \textbf{(2)} a public research artifact that implements the intervention end-to-end and is open-sourced alongside this work, and \textbf{(3)} a field-study design for evaluating how tone interacts with substantive fact-checking to shape affective polarization, engagement and the character of downstream discussion. The study is currently deployed on X; this report presents the framing, system, and protocol, with empirical results forthcoming.